\subsection{Case of Hilbert spaces}

In this section, E denotes a Hilbert space over C. We denote by \(\pi\) the canonical homomorphism of \(\mathcal{L}(\mathrm{E})\) onto the Calkin algebra \(\mathcal{Calk}(\mathrm{E})\) .

PROPOSITION 7. --- Let \(u \in \mathcal{L}(\mathrm{E})\) .

\begin{enumerate}
\def\labelenumi{\alph{enumi})}
\item
  If u is normal, then u is a Riesz endomorphism if and only if u is a Fredholm endomorphism of index 0.
\item
  If u is Hermitian, then u is a Fredholm endomorphism if and only if u is a Riesz endomorphism.
\end{enumerate}

Every Riesz endomorphism is a Fredholm endomorphism of index 0 (Proposition 5 of III, p. 46). Conversely, suppose that u is a Fredholm endomorphism of index 0, and that u is normal. Its nilspace then coincides with its kernel (EVT, V, p. 43, prop. 8), and is therefore finite-dimensional. It then follows from Lemma 2 of III, p. 45 and from the definition that u is a Riesz endomorphism.

Let us prove \(b\) ). According to \(a\) ), it suffices to check that the index of a Hermitian Fredholm endomorphism \(u\) is zero. But the orthogonal complement of the image of \(u\) (which is closed) is then equal to the kernel of \(u\) (EVT, V, p. 41, prop. 4), whence the assertion.

COROLLARY. --- Let \(u \in \mathcal{L}(\mathrm{E})\) . If u is normal, then \(\mathrm{Sp}_{0}(u)\) is empty, and if u is Hermitian, then \(\mathrm{Sp}_{\mathrm{e}}(u)\) coincides with the spectrum of \(\pi(u)\) relative to the algebra \(\mathcal{Calk}(\mathrm{E})\) .

Both assertions follow from the proposition and the definitions of the sets \(\mathrm{Sp}_{n}(u)\) for \(n \in \overline{Z}\) and of \(\mathrm{Sp}_{\omega}(u)\) , which form a partition of the essential spectrum of u (def. 2 of III, p. 83).

THEOREM 3 (Weyl). --- Let \(u \in \mathcal{L}(\mathrm{E})\) be a normal endomorphism of E. The essential spectrum of \(u\) is the intersection of the sets \(\mathrm{Sp}(u + h)\) for \(h\) ranging over \(\mathcal{L}^{\mathrm{c}}(\mathrm{E})\) .

Let \(h \in \mathcal{L}^c(\mathrm{E})\) . Since \(\mathrm{Sp}_0(u)\) is empty (corollary above), Theorem 2 of III, p. 89 implies that \(\mathrm{Sp_e}(u + h) = \mathrm{Sp_e}(u)\) . The intersection of the sets \(\mathrm{Sp}(u + h)\) therefore contains \(\mathrm{Sp_e}(u)\) .

Let \(\lambda \in \mathrm{Sp}_{\mathrm{s}}(u)\) . Denote by \(\mathrm{E}_{\lambda}\) the eigenspace of \(u\) relative to \(\lambda\) , and \(\mathrm{F}_{\lambda}\) the image of the spectral projector associated with \(u\) and \(\mathbf{C} = \{\lambda\}\) . The space E is the topological direct sum of \(\mathrm{E}_{\lambda}\) and \(\mathrm{F}_{\lambda}\) . Let \(h\) be the finite-rank endomorphism of E that vanishes on \(\mathrm{F}_{\lambda}\) and coincides on \(\mathrm{E}_{\lambda}\) with the identity. The endomorphism \(u + h\) is invertible, hence \(\lambda \notin \mathrm{Sp}(u + h)\) . The theorem follows from this.

Thus, if u is a normal endomorphism of E, and if \(h \in \mathcal{L}^{c}(\mathrm{E})\) , the spectrum of \(u + h\) can differ from the spectrum of u only by isolated points that have finite spectral multiplicity.
% semantic-fragments: legacy-input-seam
\relax{}
